\section{Implementation in an array language}
\label{sec:implementation}

The derivation is indifferent to how the factors are stored and to whether a
constraint normal is dense or a single scaled unit vector. An implementation is
not, and the difference is large enough to change which method one should use.
This section collects the three decisions that matter and measures each. They
share a moral, stated at the end of the paper: in a compiled language, layout is
chosen for cache behaviour; in an array language it additionally decides
\emph{which kernels are reachable at all}.

\subsection{What changes when a scalar operation costs a microsecond}

The reference implementations of this method --- Powell's, Turlach's Fortran, and
the C descended from it --- are written for a machine model in which a scalar
operation is nearly free and the unit of composition is the loop. An array
language inverts that. An operation reached through the interpreter costs on the
order of a microsecond whether it touches ten elements or ten thousand, so the
design question stops being \emph{how many floating-point operations?} and becomes
\emph{how many array operations, and does each reach a tuned kernel?}

Two consequences follow, and both are visible in the measurements elsewhere in
this paper. Any part of the method that is irreducibly sequential in the number of
\emph{iterations}, rather than in the number of elements, stays interpreted and
dominates at small $n$: that is the flat region below a few hundred variables in
Figure~\ref{fig:sweep}, and it is why the fast path of Section~\ref{sec:pdas},
which attacks the iteration count rather than the cost of an iteration, wins more
at these sizes than any amount of kernel tuning (Section~\ref{ssec:compare}).
Conversely, above that region the arithmetic dominates and the question is
entirely which kernel the data can be handed to.

\subsection{Storage decides which kernel is reachable}
\label{ssec:packed}

The triangular factor $R$ is held as packed columns: entry $(i,j)$ with $i \le j$
at flat offset $j(j+1)/2 + i$, so column $j$ is one contiguous run of $j+1$
values. The obvious reading is that this halves memory, and on that reading it is
an optimisation one might reasonably discard for a dense $(r,r)$ array indexed the
ordinary way. That reading is wrong, and the reason is worth measuring.

The solver reads the \emph{leading} $k \times k$ block of $R$ once per iteration,
to solve $R r = d_1$ for the dual direction~\eqref{eq:dzr}. In packed storage that
block is the leading $k(k+1)/2$ entries --- contiguous at every $k$, by
construction --- and a packed triangular solve reads it in place. In a dense
$(r,r)$ array the same block is a strided submatrix whose column stride is $r$,
not $k$; the only triangular solve that accepts it is the general LAPACK routine.

Figure~\ref{fig:layout} times exactly that one operation at the three call shapes,
with the enclosing array $\qpLayStride \times$ the working set, which is what a
solver at half its constraint budget sees.

\begin{figure}[ht]
\centering
\includegraphics[width=0.62\textwidth]{graphs/quadprog_layout.pdf}
\caption{One triangular solve of order $k$, at the three shapes the solver could
present it in. The flop count is identical in all three.}
\label{fig:layout}
\end{figure}

At $k = \qpLayKhi{}$ the packed solve takes $\qpLayPacked\,\mu$s against
$\qpLayStrided\,\mu$s for the strided view, a factor of $\qpLayRatio$, and the
ratio grows with $k$ across the range $\{\qpLayKs\}$. None of it is arithmetic:
the three computations are identical and are checked against each other.

\begin{remark}[The mechanism is not the one usually given]\label{rem:notcopy}
The natural explanation is that the strided view forces a copy at the library
boundary, and that the copy is the cost. On the evidence here that explanation is
wrong. The middle curve of Figure~\ref{fig:layout} hands the same routine a
\emph{contiguous} copy made in advance, removing any boundary copy from the timed
region --- and it is indistinguishable from the strided case, $\qpLayCopy\,\mu$s
against $\qpLayStrided\,\mu$s, a factor of $\qpLayRatioCopy$. The cost is in the
routine, not in the striding.

What packed storage buys is therefore not the avoidance of a copy but access to a
different kernel: a packed triangular solve is a level-2 BLAS operation reading
its argument in place, and no dense layout can reach it. That is a stronger claim
than the copy story and a simpler one. We record the correction because the copy
explanation is plausible, is what we assumed, and survives casual measurement ---
timing only the strided and packed shapes reproduces the headline ratio and
appears to confirm it.
\end{remark}

The price is paid in the deletion of Section~\ref{ssec:delete}, which mixes two
\emph{rows} of $R$ across a range of columns. Column offsets grow with the column
index, so a row is a strided gather in the packed layout where it would be a slice
in the dense one. Insertion is unaffected: a column is one contiguous run either
way. The asymmetry is deliberate, and it is the right way round --- the triangular
solve runs once per iteration unconditionally, the deletion chase only when a
constraint leaves.

\subsection{A bound constraint is a single nonzero}
\label{ssec:structure}

A bound $x_\rho \ge \ell$ is a column of $C$ holding one nonzero, $c_i = \nu
e_\rho$, and a box-constrained problem consists of nothing else. Three of the
per-iteration quantities then degenerate to indexing:
\begin{equation}
  c_i\T x = \nu\, x_\rho,
  \qquad
  d = J\T c_i = \nu\, (J_{\rho\cdot})\T,
  \qquad
  c_i\T z = \nu\, z_\rho ,
  \label{eq:unit}
\end{equation}
replacing an $O(nm)$ product, an $O(n^2)$ product and an $O(n)$ dot product
respectively. The identities are exact in floating point --- each is a single
multiplication where the general form is a sum of $n$ terms of which $n-1$ are
zero, and adding exact zeros in IEEE~754 changes nothing --- so the fast path
computes the same values, not merely close ones.

Two details of the detection matter more than the arithmetic.

It is performed \emph{per column}, not per matrix. The case worth optimising is
mixed: a long-only portfolio problem carries one dense budget row
$\mathbf{1}\T x = 1$ alongside $2n$ bounds, and an all-or-nothing test would
observe that single dense column and send the entire problem down the general
path. Per-column detection keeps the $2n$ bounds cheap and pays the general cost
only on the one column that needs it.

And an all-zero column must not be mistaken for a unit column with $\nu = 0$. It
expresses $0 \ge b_i$, the case excluded in Proposition~\ref{prop:inner}; treating
it as a scaled unit vector would scale a row of $J$ by zero and lose the
infeasibility verdict that Section~\ref{sec:termination} owes it.

Measured end to end over $n \in \{\qpStructNs\}$, by comparing a box problem
against the same problem with each column perturbed so that no column is a single
nonzero, detection is worth up to $\qpStructRatio\times$, at
$n = \qpStructN{}$, with both variants returning the same minimiser exactly. The
gain is modest at the sizes where a solve is dispatch-bound and grows with $n$ as
the slack product starts to matter, which is the expected shape: the product runs
once per outer iteration over all $m$ constraints, and $m = 2n$ on a box.

\subsection{The BLAS thread count, which is not the solver's to choose}
\label{ssec:threads}

Because the method pushes its work into BLAS calls, its performance inherits
whatever threading decision the installed library makes, and that decision is not
uniformly benign. On one configuration --- Linux, NumPy built against OpenBLAS,
and a thread count above the physical core count --- oversubscription is
catastrophic rather than merely wasteful, because at these sizes a level-2 kernel
has too little work per call to amortise a many-way spin-waiting barrier, and
under simultaneous multithreading the spinning threads contend with the working
ones for the same physical core. The collapse is not gradual: it appears when the
library crosses its internal threshold for threading a given kernel, so a run can
look healthy at one size and be an order of magnitude slower at twice that size.

We report this rather than measure it here, and the distinction matters. The
machine these experiments were run on uses a BLAS that exposes no thread control
at all (Section~\ref{sec:reproducibility}), so the effect is not reproducible from
this paper's scripts; the sweeps behind it are distributed with the
package~\cite{cvxquadprog} and are the appropriate reference. What belongs in a
paper about the method is the structural observation, which survives any
particular library: a solver whose cost is concentrated in short level-2 calls is
exposed to its BLAS's threading defaults in a way that a solver dominated by large
level-3 calls is not, and the exposure is asymmetric --- capping threads costs
little when it is unnecessary and saves an order of magnitude when it is not.
The implementation therefore caps the count on that one configuration and leaves
every other alone, on the principle that a default should be chosen to bound the
worst case rather than to optimise the common one.
